%--- iliad ---
% summary: >-
%   What must a next-token predictor remember about the past? Causal states,
%   belief states and their geometry, from HMMs to generalised HMMs, ending
%   with the evidence that transformers learn this structure.
%--- end ---
\documentclass[11pt]{article}
\IfFileExists{iliad.sty}{\usepackage[boxes]{iliad}}{\usepackage[boxes]{../iliad}}
% float: the Mess3 solution carries a figure inside its box, where a float
% cannot go; [H] sets it in place.
\usepackage{float}

\title{Computational Mechanics}
\author{\authorname{Xavier Poncini}\\ \affiliation{Simplex}
  \and \authorname{Adam Shai}\\ \affiliation{Simplex}
  \and \authorname{Paul Riechers}\\ \affiliation{Simplex}}
\date{}

\begin{document}
\maketitle

\section{Prerequisites}\label{prerequisites}

Familiarity with the goals and basic methodology of mechanistic interpretability, 
including the notions of features and circuits; the linear representation hypothesis; 
linear probes as a method for reading off internal representations; and a basic 
understanding of the transformer architecture. 

Students should be comfortable with the idea that one can train a linear map 
from model activations to some target structure and evaluate its 
quality (e.g. via MSE or \(R^2\)).

\textbf{Mathematical background.} The module is mathematically self-contained: 
all formal definitions (HMMs, GHMMs, belief states, MSPs) are introduced from scratch. 
However, students will engage with the material more fluently if they are comfortable with the following:

\begin{itemize}
\item
  \textbf{Linear algebra:} row-stochastic matrices, rank and invertibility, null spaces (kernels).
\item
  \textbf{Probability:} conditional probability, Bayes\textquotesingle{} rule, probability distributions over finite sets, the probability simplex.
\item
  \textbf{Basic machine learning:} next-token prediction, loss functions (cross-entropy), the concept of model activations.
\end{itemize}

\begin{learningoutcomes}
\begin{itemize}
\item
  Construct their own hidden Markov models.
\item
  Explain the motivation for introducing generalised hidden Markov models (GHMMs).
\item
  Explain why the belief state is a useful object for making predictions about GHMM data.
\item
  Compute the belief state of a GHMM.
\item
  Explain the difference between the belief state of a HMM and the belief state of a GHMM (that is itself not an HMM).
\item
  Explain the relationship between belief states and next-token observation probabilities for GHMMs.
\item
  Compute the mixed state presentation (MSP) for GHMMs.
\item
  Explain why generating data is easier than predicting data.
\item
  Explain the evidence for the claim: transformers represent belief state geometry in their residual stream.
\item
  Develop hypotheses for where transformers represent: belief states, observation probabilities, log observation probabilities.
\end{itemize}
\end{learningoutcomes}

\section{Content}\label{content}

\subsection{Fast-track}\label{fast-track}

Read \href{https://arxiv.org/pdf/2405.15943}{Transformers represent belief state geometry in their residual stream}. Focus on being able to answer the following questions:

\begin{itemize}
\item
  What is a hidden Markov model (HMM)?
\item
  What is a belief state?
\item
  What is the mixed state presentation (MSP)?
\item
  About the map learnt from activations to belief states:
  \begin{itemize}
  \item
    What is the source?
  \item
    What is the target?
  \item
    How is the quality of the map evaluated?
  \end{itemize}
\end{itemize}

% The roadmap is the Schedule from the source repo's README
% (iliad-team/iliad-intensive-C.3), which is what the day was taught from on
% 10 Sep 2026. Its links pointed at the repo's own PDFs; they point at this
% page's decks and sections instead. Words are the README's.
\subsection{Schedule}\label{schedule}

\begin{itemize}
\item
  \textbf{10:00--10:30} --- \href{https://iliad-intensive.org/downloads/computational-mechanics/computational-mechanics-slides.pdf}{Lecture: Overview and scope}
\item
  \textbf{10:30--11:15} --- \href{https://iliad-intensive.org/downloads/computational-mechanics/computational-mechanics-slides-2-predicting-the-future.pdf}{Lecture: Predicting the future}
\item
  \textbf{11:15--11:30} --- \emph{Break}
\item
  \textbf{11:30--12:30} --- \hyperref[sec:predicting-the-future]{{Exercises}: Predicting the future}
\item
  \textbf{12:30--13:30} --- \emph{Lunch}
\item
  \textbf{13:30--14:15} --- \href{https://iliad-intensive.org/downloads/computational-mechanics/computational-mechanics-slides-3-representing-the-past.pdf}{Lecture: Representing the past}
\item
  \textbf{14:15--15:45} --- \hyperref[sec:representing-the-past]{{Exercises}: Representing the past}
\item
  \textbf{15:45--16:00} --- \emph{Break}
\item
  \textbf{16:00--17:45} --- \hyperref[sec:belief-geometry]{Reading \& Discussion: Transformers represent belief geometry}
\item
  \textbf{17:45--18:00} --- Feedback
\end{itemize}

% Parts 2-4 are the source repo's written-exercises/CompMechExercisesPt{2,3}
% and reading-guides/CompMechReadingGuidePt4, run through a scripted
% transform (environment names only; every word and formula is the
% author's). Part 1 is lecture-only: its deck is the first Slides row.
\section{Predicting the future}\label{sec:predicting-the-future}
\input{part2}

\section{Representing the past}\label{sec:representing-the-past}
\input{part3}

\section{Identifying belief geometry in transformers}\label{sec:belief-geometry}
\input{part4}

% From the source repo's README, "Python exercises: HMM from scratch".
\section{Python exercises: HMM from scratch}\label{sec:python-exercises}

The Python exercises are ARENA-style notebooks in two parts. Each part has an
\textbf{exercises} notebook (with \texttt{\# YOUR CODE HERE} stubs, inline tests, and
collapsible solutions) and a fully worked \textbf{solutions} notebook.

\begin{itemize}
\item
  \textbf{Part 1 --- Sequence probabilities \& next-token distributions} (exercises) --- \href{https://colab.research.google.com/github/iliad-team/iliad-intensive-C.3/blob/main/python-exercises/colab/part1_sequence_probabilities_exercises_colab.ipynb}{Open in Colab}
\item
  Part 1 (solutions) --- \href{https://colab.research.google.com/github/iliad-team/iliad-intensive-C.3/blob/main/python-exercises/colab/part1_sequence_probabilities_solutions_colab.ipynb}{Open in Colab}
\item
  \textbf{Part 2 --- Belief states} (exercises) --- \href{https://colab.research.google.com/github/iliad-team/iliad-intensive-C.3/blob/main/python-exercises/colab/part2_belief_states_exercises_colab.ipynb}{Open in Colab}
\item
  Part 2 (solutions) --- \href{https://colab.research.google.com/github/iliad-team/iliad-intensive-C.3/blob/main/python-exercises/colab/part2_belief_states_solutions_colab.ipynb}{Open in Colab}
\end{itemize}

The Colab notebooks are self-contained: their setup cells write the helper
modules into the session, so there is nothing else to install or download.

\section{Learn more}\label{learn-more}

\begin{itemize}
\item
  Beyond toy architectures -- in-context learning
  \begin{itemize}
  \item
    LLMs have an amazing ability to adapt internal representations in-context -- {[}Read: ICLR: \href{https://arxiv.org/pdf/2501.00070}{In-Context Learning of Representations}{]}
  \item
    LLMs can predict HMM data in-context -- {[}Read: \href{https://arxiv.org/pdf/2506.07298}{Pre-trained Large Language Models Learn Hidden Markov Models In-context}{]}
  \item
    This suggests that internal representations of LLMs predicting GHMM data in-context, may resemble belief states.
  \end{itemize}
\item
  Processes consisting of GHMMs with richer structure {[}Read: \href{https://arxiv.org/pdf/2602.02385v1}{Transformers learn factored representations}{]}
\item
  Review article of computational mechanics {[}Read: \href{https://www.nature.com/articles/nphys2190}{Between order and chaos}{]}
\item
  \href{https://www.simplexaisafety.com/}{Simplex} is a non-profit research organisation working on applying computational mechanics to AI safety.
\end{itemize}

\end{document}
